\section{Competition and information}\label{sec:payoffs}

\subsection{Acquisition profits and target value}\label{sec:profits}

The auction determines both the return to preparation and the claim that informed investors trade. Consider a
prepared challenger with value \(\theta\ge p\). If it wins, it pays \(\max\{p,R\}\); if it loses, the
incumbent pays \(\theta\). The target therefore receives \(\max\{p,\min(R,\theta)\}\), and the challenger
keeps \((\theta-\max\{p,R\})_+\). A prepared challenger with \(\theta<p\) cannot bid. The target then receives
the reserve if the incumbent's value meets it, exactly as without entry.

For every value at or above the reserve, the two expected amounts add up to the value:
\[
t(\theta)+g(\theta)=\theta,\qquad \theta\ge p.
\]
The second-price rule passes the whole value of an admissible challenger to the target and the challenger
together. An admissible value also raises proceeds to at least the reserve, so \(t_j\ge t_0\) for both
classes.

In the point-mass benchmark with the uniform incumbent, integration over \(R\) gives
\begin{equation}\label{eq:payoffs}
\begin{aligned}
t_0&=p\Bigl(1-\frac pr\Bigr),&
t_H&=\frac r2+\frac{p^2}{2r},\\
t_L&=\ell-\frac{\ell^2-p^2}{2r},&
g_H&=h-\frac r2-\frac{p^2}{2r},\\
g_L&=\frac{\ell^2-p^2}{2r},&
\DT(r)&=\frac{(r-\ell)^2}{2r}.
\end{aligned}
\end{equation}
The relevant derivatives are
\begin{equation}\label{eq:derivatives}
\DT'(r)=\frac12-\frac{\ell^2}{2r^2}>0,\qquad
g_H'(r)=-\frac12+\frac{p^2}{2r^2}<0,\qquad
g_L'(r)=-\frac{\ell^2-p^2}{2r^2}<0.
\end{equation}

Thus a stronger incumbent lowers acquisition profit at every fixed belief. At the same time, it makes target
proceeds more sensitive to challenger quality. Proposition~\ref{prop:returns} extends this opposition to every
continuous incumbent distribution and every pair of classes that satisfies Assumption~\ref{ass:classes}.

\begin{proposition}[Competition and the two returns to information; analytical]\label{prop:returns}
Let the incumbent's value have a continuous distribution \(F\) on \([0,\bar r]\), extended by unity above
\(\bar r\). For any classes \(G_H,G_L\) with finite means,
\begin{equation}\label{eq:spread}
\begin{aligned}
\DT(F)&=pF(p)D(p^-)+\int_p^\infty[1-F(y)]\,D(y)\,dy,\\
g_j(F)&=\int_p^\infty[1-G_j(y)]\,F(y)\,dy,\qquad j\in\{H,L\}.
\end{aligned}
\end{equation}
Let \(\tilde F\le F\) be a continuous first-order stochastic strengthening of \(F\).
\begin{enumerate}[label=(\alph*),leftmargin=2em]
\item For every pair of classes, the strengthening weakly decreases \(g_H\) and \(g_L\), and hence gross
acquisition profit at every fixed belief.
\item Under Assumption~\ref{ass:classes}, the strengthening weakly increases the target-payoff spread.
\end{enumerate}
The decrease in \(g_j\) is strict when \(\int_p^\infty[1-G_j(y)][F(y)-\tilde F(y)]\,dy>0\). The increase in
(b) is strict when \(\int_p^\infty[F(y)-\tilde F(y)]\,D(y)\,dy>0\).
\end{proposition}

The proof works one incumbent value at a time. At a value \(R\) that meets the reserve, the target receives
\(\min\{R,\max\{\theta,p\}\}\). The expected class difference in proceeds is then \(\int_p^RD(y)\,dy\). Read
the integrand as slices \(D(y)\,dy\) of the class value gap at each level \(y>p\). The slice at \(y\) reaches
the target when the incumbent outbids the challenger at that level and pays its bid. It stays with the
challenger otherwise.

A stronger incumbent moves slices to the target. The spread rises and acquisition profit falls. The reserve
adds one term. Below the reserve, only an admissible challenger buys, at \(p\), and the class difference is
\(pD(p^-)\). A stronger incumbent fails the reserve less often, so that term shrinks. It can outweigh the moved
slices when the low class is more often inadmissible.

Under Assumption~\ref{ass:classes}, every value meets the reserve, and the term is zero.
Appendix~\ref{app:payoffs} gives the proof. It also gives, in Proposition~1(c), the weakest restriction under
which the spread rises for every strengthening.

The sum of the two lines of \eqref{eq:spread} under Assumption~\ref{ass:classes} is
\[
\DT(F)+g_H(F)-g_L(F)=\E_H[\max\{\theta,p\}]-\E_L[\max\{\theta,p\}].
\]
The identity holds for every incumbent distribution. In the benchmark, the right side is \(h-\ell\) at every
\(r\). A stronger incumbent leaves the value gap between the classes unchanged. It moves the gap from the
challenger's profit to target proceeds.

The spread approaches zero as \(r\) approaches \(\ell\), since the incumbent then almost never outbids a
low-value challenger. As \(r\) increases, the spread rises while gross acquisition profit falls at each belief.
Whether entry rises depends on how this change in payoffs affects the information in the price.

\subsection{What the challenger learns from the price}\label{sec:inference}

The challenger sees one price. This section states what that price can reveal under every investor strategy, and
what it reveals in equilibrium. Write \(\delta_q\) for the point mass at the order \(q\), and
\(b_0=(1-\lambda)(1-2\beta)\) for the probability of a liquidity abstention. Under any order distributions
\(\sigma_H,\sigma_L\), Bayes' rule in odds form reads
\begin{equation}\label{eq:odds}
\frac{\mu(x)}{1-\mu(x)}=\frac{\pi}{1-\pi}\cdot\frac{f_H(x)}{f_L(x)}.
\end{equation}

\begin{lemma}[Order probabilities and beliefs; analytical]\label{lem:beliefs}
Fix any order distributions \(\sigma_H\) and \(\sigma_L\).
\begin{enumerate}[label=(\alph*),leftmargin=2em]
\item \emph{Full support.} \(f_j(x)\ge(1-\lambda)\eta(x)>0\) for every order \(x\) and class \(j\).
\item \emph{Nonzero orders.} For \(x\in\{-1,1\}\), \(b\le f_j(x)\le b+\lambda\). Hence
\(m(\pi)\le\mu(x)\le M(\pi)\).
\item \emph{The zero order.} \(b_0\le f_j(0)\le b_0+\lambda\), and \(b_0\ge b\). Hence
\(m(\pi)\le\mu(0)\le M(\pi)\).
\item \emph{Full orders.} If \(\sigma_H=\delta_1\) and \(\sigma_L=\delta_{-1}\), the posteriors after the orders
\((-1,0,1)\) are \(\bigl(m(\pi),\pi,M(\pi)\bigr)\).
\end{enumerate}
\end{lemma}

Each class reaches an order in a given direction through a liquidity order with probability \(b\), and the
informed investor adds at most \(\lambda\). The likelihood ratio \(f_H(x)/f_L(x)\) therefore lies in
\([1/(1+u),1+u]\), and by \eqref{eq:odds} the odds of \(\mu(x)\) lie between the odds of \(m(\pi)\) and those of
\(M(\pi)\). At the zero order the liquidity mass is \(b_0\), and \(b_0\ge b\) is the inequality \(\beta\le1/3\).
With \(\beta>1/3\), a profile in which class \(H\) abstains and class \(L\) never abstains gives \(\mu(0)>M(\pi)\).
Appendix~\ref{app:beliefs} gives the proofs of this section.

Part~(c) is the only use of the bound \(\beta\le1/3\), and Corollary~\ref{cor:exclusion} uses it.
Propositions~\ref{prop:ccc} and~\ref{prop:coexistence} and Corollary~\ref{prop:priors} use parts (a), (b) and (d)
only, so they hold for every \(\beta<1/2\).

The challenger's belief at a price is an average of the posteriors of the orders with that price, so it obeys the
same bounds. A preparation decision that needs a belief above \(M(\pi)\) therefore follows no price.

\begin{corollary}[High-cost preparation under every profile; analytical]\label{cor:exclusion}
Let \(g_H\ge g_L\). If \(B(M(\pi))<c_H\), then no profile that satisfies (E1)--(E3), in or out of equilibrium,
has high-cost preparation after any order.
\end{corollary}

Corollary~\ref{cor:exclusion} is the ceiling of the model. Against a strong enough incumbent, gross profit falls
below the high cost even at the top posterior, and no order can recruit an expensive challenger.

When \(g_H>g_L\), write
\[
\tau=\frac{c_H-g_L}{g_H-g_L}
\]
for the preparation threshold. Gross profit covers the high cost exactly when the belief reaches the threshold:
\(B(\mu)\ge c_H\) if and only if \(\mu\ge\tau\). For the uniform incumbent I write \(\tau(r)\).

\begin{lemma}[Entry and prices; analytical]\label{lem:prices}
Let \(\DT>0\) and \(g_H>g_L\). Fix a tuple that satisfies (E1)--(E4) and (E6).
\begin{enumerate}[label=(\alph*),leftmargin=2em]
\item \emph{Floor.} \(e(x)\ge\rho\) for every order \(x\).
\item \emph{Price.} \(P(x)=t_0+e(x)\bigl[t_L-t_0+\DT\,\mu(x)\bigr]\) for every order \(x\).
\item \emph{Sufficiency.} If \(P(x)=P(x')\), then \(\mu(x)=\mu(x')\). The challenger's belief at the price \(P(x)\)
is \(\mu(x)\).
\item \emph{Entry.} \(e(x)=\rho\) if \(\mu(x)<\tau\), and \(e(x)=1\) if \(\mu(x)\ge\tau\).
\item \emph{Monotonicity.} If \(\mu(x)>\mu(x')\), then \(e(x)\ge e(x')\) and
\(P(x)-P(x')\ge\rho\DT\,[\mu(x)-\mu(x')]>0\).
\end{enumerate}
\end{lemma}

Part~(c) is the step from orders to prices. Suppose that two orders have one price. The challenger cannot tell them
apart, so entry after them is the same, and it is at least \(\rho\). The price formula then gives the same posterior
after both orders. The challenger can therefore read the market makers' posterior from the price, and its gross
profit at the price \(P(x)\) is \(B(\mu(x))\). Under rule CM, part~(d) at \(\mu(x)=\tau\) allows any
\(e(x)\in[\rho,1]\); the other parts hold under both challenger rules.

The zero-cost challenger does three jobs in the proofs. Entry stays at least \(\rho\), so the price reveals the
posterior. The investor's return after every order stays at least a fixed share of the spread
(Lemma~\ref{lem:trading} below). Zero-cost preparation does not respond to the price, so information can only add
entry.

\subsection{Informed trading}\label{sec:informed}

Rational pricing removes the anticipated increase in proceeds from the investor's advantage. The next lemma states
what remains, under every candidate tuple.

\begin{lemma}[Informed trading; analytical]\label{lem:trading}
Let \(\DT\ge0\). Fix a tuple that satisfies (E1) and (E4), with entry \(e(x)\in[\rho,1]\) after every order.
Conditional on its arrival, the investor in each class earns
\[
\begin{aligned}
U_H(1)&=e(1)\DT[1-\mu(1)]-k,&\quad U_H(0)&=0,&\quad U_H(-1)&=-e(-1)\DT[1-\mu(-1)]-k,\\
U_L(-1)&=e(-1)\DT\,\mu(-1)-k,&\quad U_L(0)&=0,&\quad U_L(1)&=-e(1)\DT\,\mu(1)-k.
\end{aligned}
\]
Write \(K_H=e(1)\DT[1-\mu(1)]\) and \(K_L=e(-1)\DT\,\mu(-1)\) for the gross returns of the two correct orders.
\begin{enumerate}[label=(\alph*),leftmargin=2em]
\item A wrong-signed order earns at most \(-k\), so abstention strictly beats it.
\item Class \(H\) buys if \(K_H>k\) and abstains if \(K_H<k\). At \(K_H=k\), the purchase and abstention are its
two best orders. Class \(L\) sells if \(K_L>k\) and abstains if \(K_L<k\). At \(K_L=k\), the sale and abstention
are its two best orders.
\item Under every candidate profile, \(\rho\DT[1-M(\pi)]\le K_H\le\DT\) and \(\rho\DT\,m(\pi)\le K_L\le\DT\).
Hence \(\rho\,\kappa(\pi)\DT\le K_j\le\DT\), where \(\kappa(\pi)=\min\{m(\pi),1-M(\pi)\}\).
\item \(K_H\) depends on the candidate only through \(\mu(1)\) and \(e(1)\), and \(K_L\) only through
\(\mu(-1)\) and \(e(-1)\).
\end{enumerate}
\end{lemma}

The investor's gross return is the entry probability times the spread times the market's residual uncertainty
about the class. The market makers price the entry response, but a liquidity order in the same direction is always
possible, so some uncertainty survives every order. The lower bound in part~(c) uses the entry floor \(\rho\) and
the posterior bounds of Lemma~\ref{lem:beliefs}. The upper bound is the spread. Both bounds hold under every
candidate profile, mixed or pure, and their comparison with \(k\) decides when information survives in
equilibrium.

The bound \(\kappa(\pi)\) has a simple form. Since \(1-M(\pi)=m(1-\pi)\), \(\kappa(\pi)=m(\min\{\pi,1-\pi\})\). It
equals \(m(1/2)=1/(2+u)\) at the prior \(1/2\) and is smaller at every other prior. The class that the prior
favors binds, because the price already leans toward it.

The no-entry proceeds \(t_0\) cancel from each return, because they move the price and the terminal payoff
together. The terms that matter are the entry gains \(t_H-t_0\) and \(t_L-t_0\), which are nonnegative.

Part~(d) splits the investor's problem. Entry after the zero order enters no payoff. Each class therefore compares
one gross return with the trading cost: class \(H\) the return of a purchase, class \(L\) the return of a sale.
Section~\ref{sec:coexistence} uses this split to list every equilibrium along the benchmark path.
